% ============================================================
%  Chapter 8: Evaluation — Correctness, Security, and Empirical Framework
%  Formal proofs, threat model, benchmark and user study design
% ============================================================

\chapter{Evaluation: Correctness, Security, and Empirical Evidence}
\label{ch:evaluation}

\section{Formal Correctness Proofs of the Admission Pipeline}

The \texttt{star-toml} implementation must satisfy three core correctness properties:
(1) every admitted configuration is valid;
(2) the witness is deterministic and collision-resistant;
(3) the canonical form is idempotent under re-serialization.

\begin{theorem}[Admission Validity]
\label{thm:admission-validity}
For all configurations $x \in \CC^{\mathrm{complete}}$ and bounded artifact space $\mathcal{X}_B$:
\[
  x \in \mathcal{A} \implies
  \left(
    \forall k \in \Sigma_{\mathrm{req}}, x_k \neq \bot
    \land \forall k \in \Sigma, \phi_k(x_k) = 1
  \right)
\]
In other words: every admitted configuration has all required fields defined
and satisfies all declared field constraints.
\end{theorem}

\begin{proof}
$x \in \mathcal{A}$ iff $\qc(x, W, \sigma, w, \EE) = 1$ (Definition~\ref{def:brce}).
This requires $\mathrm{BS}(x, \Phi) = 1$ (Definition~\ref{def:bs}).
BS is defined as $\bigwedge_{k \in \Sigma} \phi_k(x_k) = 1$ (Definition~\ref{def:bs}),
which enforces all field constraints.
Additionally, the stage function $f_4$ (Deserialization) succeeds only if
all required fields are defined in the merged configuration $x^*$
(checked by stage $f_2$, TransformLawful).
Therefore every admitted $x$ satisfies the claimed properties. \qed
\end{proof}

\begin{theorem}[Canonical Serialization Injectivity]
\label{thm:canonical-injectivity}
The canonical serialization $\kappa : \CC^{\mathrm{complete}} \to \Sigma^*$
is injective. For all $x, y \in \CC^{\mathrm{complete}}$:
\[
  \kappa(x) = \kappa(y) \implies x = y
\]
\end{theorem}

\begin{proof}
By definition (Definition~\ref{def:witness}), $\kappa$ is a total, deterministic,
fixed-order serialization with no formatting variation.
Two distinct configurations $x \neq y$ produce distinct field sequences;
the fixed-order, no-variation serialization produces distinct byte strings.
(If $\kappa(x) = \kappa(y)$, then the field sequences are identical,
hence $x = y$.) \qed
\end{proof}

\begin{theorem}[Witness Determinism]
\label{thm:witness-determinism}
For all configurations $x$, layer orderings $\ell = (\sigma, \{x^{(i)}\})$,
and environment bindings $e$, the witness function is deterministic:
\[
  \omega(x, \ell, e) = \omega(x', \ell', e')
  \text{ whenever } (x, \ell, e) = (x', \ell', e')
\]
\end{theorem}

\begin{proof}
$\omega$ is defined as:
\[
  \omega(x, \ell, e) = \mathrm{BLAKE3}\!\left(
    \bigoplus_{i=1}^{n} \mathrm{BLAKE3}(x^{(i)})
    \;\|\; \mathrm{BLAKE3}(\mathrm{enc}(\sigma))
    \;\|\; \mathrm{BLAKE3}(\mathrm{enc}(e))
    \;\|\; \mathrm{BLAKE3}(\kappa(x))
  \right)
\]
Each component (BLAKE3, XOR, concatenation, canonical serialization) is
deterministic on equal inputs. Therefore $\omega$ is deterministic. \qed
\end{proof}

\begin{theorem}[Witness Uniqueness over Variable-Order Configurations]
\label{thm:witness-order-independence}
If two configurations $x_1, x_2$ are equal as mappings (same field values)
but differ only in the order of variable-environment substitutions,
and both have the same sorted layer ordering $\sigma$ (after sort-normalization),
then their witnesses are equal:
\[
  \mathrm{same\_fields}(x_1, x_2)
  \land \sigma(x_1) = \mathrm{sort}(\sigma(x_1)) = \sigma(x_2)
  \implies \omega(x_1, \ell_1, e_1) = \omega(x_2, \ell_2, e_2)
\]
\end{theorem}

\begin{proof}
The witness function includes $\mathrm{BLAKE3}(\kappa(x))$, where $\kappa$ is
the canonical serialization with fixed field ordering (Definition~\ref{def:witness}).
If $x_1$ and $x_2$ have the same field values, then $\kappa(x_1) = \kappa(x_2)$.
The layer ordering $\sigma$ is included as $\mathrm{BLAKE3}(\mathrm{enc}(\sigma))$.
If both are sort-normalized to the same $\sigma$, they contribute identically
to the witness. Environment substitutions affect the final $e$ binding, but
if done in any order over the same variable-value mappings, $e$ is the same.
Therefore all components of $\omega$ are identical, so the witnesses are equal. \qed
\end{proof}

\begin{corollary}[Save Canonical Succeeds on Validated Configurations]
For all $x \in \mathrm{Validated}\langle T\rangle$:
\[
  \mathrm{save\_canonical}(x) \text{ succeeds and produces } \kappa(x)
\]
\end{corollary}

\begin{proof}
Validated configurations have passed the BS predicate (Definition~\ref{def:bs}),
ensuring all required fields are defined and all constraints are satisfied.
The canonical serialization $\kappa$ is a total function on complete configurations.
Therefore save_canonical always succeeds. \qed
\end{proof}

\section{Threat Model: BLAKE3 Witness Security}

The witness $\omega(x, \ell, e)$ is the cryptographic commitment to the
configuration's provenance tuple $(x, \ell, e)$. An adversary might attempt
three attack classes:

\subsection{Attack 1: Configuration Forgery}

\textbf{Threat:} An adversary forges configuration $x'$ such that
$\omega(x', \ell, e) = \omega(x, \ell, e)$ for the same provenance tuple,
and $x' \neq x$.

\textbf{Defense:}
By Theorem~\ref{thm:canonical-injectivity}, distinct configurations produce
distinct canonical forms: $\kappa(x) \neq \kappa(x')$.
The witness includes $\mathrm{BLAKE3}(\kappa(x))$ as a component.
BLAKE3 collision resistance (Aumasson et al., 2020) guarantees
$\Pr[\mathrm{BLAKE3}(A) = \mathrm{BLAKE3}(B) \mid A \neq B] \leq 2^{-128}$.
Therefore:
\[
  \Pr[\text{adversary forges } x' \text{ with same witness}] \leq 2^{-128}
\]
The work required is $2^{256}$ BLAKE3 invocations (preimage resistance).

\subsection{Attack 2: Chain Tampering}

\textbf{Threat:} An adversary modifies an earlier configuration in the fleet
(e.g., a database config deployed 6 months ago), then attempts to re-admit
it with a new witness that appears valid in an audit log.

\textbf{Defense:}
The witness includes $\bigoplus_{i=1}^{n} \mathrm{BLAKE3}(x^{(i)})$,
the rolling XOR hash of all source files.
Any edit to any source file $x^{(j)}$ changes $\mathrm{BLAKE3}(x^{(j)})$,
hence changes the XOR term, hence changes the final witness.
An adversary cannot tamper with a prior configuration and maintain its
original witness. The receipt (Definition~\ref{def:witness}) immutably
stores the prior witness.

\[
  \text{Tampering detected: } \omega_{\mathrm{recomputed}} \neq \omega_{\mathrm{stored}}
\]

\subsection{Attack 3: Cross-Domain Replay}

\textbf{Threat:} An adversary admits configuration $x$ in staging (domain A),
obtains its witness $\omega_A(x)$, then attempts to replay it in production
(domain B) by cloning the receipt.

\textbf{Defense:}
The witness includes the environment snapshot $e$ as a component
(Definition~\ref{def:witness}). The environment varies by domain:
\[
  e_{\mathrm{staging}} \neq e_{\mathrm{production}}
  \implies \omega_A(x) \neq \omega_B(x)
\]
Additionally, the genesis seed (if implemented) includes the domain identifier.
The receipt, stored immutably, certifies that the configuration was admitted
\emph{in domain A under environment $e_A$}.

Attempting to use the receipt in domain B will fail: the domain identifier
does not match, and the recomputed witness differs.

\begin{remark}[Adversary Capabilities and Limitations]
This threat model assumes:
\begin{itemize}
  \item \textbf{Attacker has file-system read access} but not write access
        to the pipeline logs or canonical storage.
  \item \textbf{Attacker does not control the BLAKE3 hash function.}
  \item \textbf{Attacker has $< 2^{128}$ computational budget.}
        (True for any realistic scenario; $2^{128}$ operations is infeasible.)
\end{itemize}
Under these assumptions, witness forgery is infeasible.
\end{remark}

\section{Benchmark Framework Design}

To evaluate the admission pipeline's operational overhead, we propose a
benchmark framework with three metrics: \emph{compilation overhead},
\emph{admission latency}, and \emph{memory footprint}.

No real data are presented here; this section structures the framework
for future empirical measurement.

\subsection{Metric 1: Compilation Overhead}

\textbf{Hypothesis:} The \texttt{\#[derive(Validate)]} macro adds negligible
compile-time overhead compared to hand-implemented validation.

\textbf{Setup:}
\begin{itemize}
  \item Generate $N \in \{10, 100, 1000\}$ TOML schemas of varying complexity.
  \item Complexity: number of fields $|\Sigma| \in \{5, 50, 500\}$;
        number of constraints $|Phi| \in \{5, 50, 500\}$.
  \item For each schema, measure:
    \begin{itemize}
      \item $T_{\mathrm{macro}}$: compilation time with \texttt{\#[derive(Validate)]}.
      \item $T_{\mathrm{hand}}$: compilation time with hand-written validator.
      \item Binary size increase.
    \end{itemize}
\end{itemize}

\textbf{Metrics:}
\[
  \mathrm{Overhead}_{\mathrm{compile}} = \frac{T_{\mathrm{macro}} - T_{\mathrm{hand}}}{T_{\mathrm{hand}}} \times 100\%
\]

\textbf{Expected Result:} $\leq 5\%$ compile-time overhead; negligible binary size impact.

\subsection{Metric 2: Admission Latency}

\textbf{Hypothesis:} The full pipeline (stages $f_1$ through $f_9$) completes
in sub-second time for typical configurations (1--10 KB), with latency
scaling linearly in configuration size.

\textbf{Setup:}
\begin{itemize}
  \item Generate synthetic configurations of sizes:
        $|x| \in \{1\text{ KB}, 10\text{ KB}, 100\text{ KB}, 1\text{ MB}\}$.
  \item Complexity levels: simple (10 fields), moderate (100 fields),
        complex (1000 fields).
  \item For each, measure:
    \begin{itemize}
      \item $\tau_1$: time for stages $f_1$--$f_4$ (parsing, merging, env resolution).
      \item $\tau_5$: time for stage $f_5$ (validation and counterexample checking).
      \item $\tau_6$--$\tau_9$: time for stages $f_6$--$f_9$ (freezing, serialization, witness).
      \item $\tau_{\mathrm{total}} = \tau_1 + \tau_5 + (\tau_6 + \cdots + \tau_9)$.
    \end{itemize}
  \item Run 100 trials per configuration; report mean and 99th percentile.
\end{itemize}

\textbf{Metrics:}
\[
  \mathrm{Latency}(|x|) = \mathbb{E}[\tau_{\mathrm{total}}]
\]
\[
  \mathrm{Linearity} = \frac{\mathrm{Latency}(1\text{ MB})}{\mathrm{Latency}(1\text{ KB})} \approx 1000
\]

\textbf{Expected Result:} $\tau_{\mathrm{total}} < 500\text{ ms}$ for 100 KB configs;
linear scaling up to 1 MB.

\subsection{Metric 3: Memory Footprint}

\textbf{Hypothesis:} The \texttt{AdmittedConfig}\langle T\rangle wrapper adds
$\leq 10\%$ memory overhead compared to the bare type $T$.

\textbf{Setup:}
\begin{itemize}
  \item For each configuration type $T$, measure:
    \begin{itemize}
      \item $M_T = \text{size of } T \text{ in memory (after deserialization)}$.
      \item $M_{\mathrm{wrapper}} = \text{size of } \texttt{AdmittedConfig}\langle T\rangle$
            (includes $T$, witness digest, and metadata).
      \item Peak memory during pipeline (max allocated across all stages).
    \end{itemize}
\end{itemize}

\textbf{Metrics:}
\[
  \mathrm{Overhead}_{\mathrm{mem}} = \frac{M_{\mathrm{wrapper}} - M_T}{M_T} \times 100\%
\]

\textbf{Expected Result:} $\leq 10\%$ overhead; peak memory $\leq 3 \times M_T$.

\section{User Study Framework: Hypothetical Developer Experiment}

To evaluate the effectiveness of the admission system in helping developers
write correct configurations, we propose a randomized controlled experiment.

\subsection{Study Design}

\textbf{Participants:} $N = 10$ software engineers with 2--10 years of experience,
split into control and treatment groups ($N/2$ each).

\textbf{Task:} Write and debug a TOML configuration file for a distributed
service with:
\begin{itemize}
  \item 30 required fields (database URL, ports, TLS paths, timeouts, etc.).
  \item 10 declared constraints (port ranges, timeout bounds, path safety).
  \item 3 layers (base, environment-specific, local overrides).
  \item Intentional bugs embedded in the setup: e.g., port 0, missing required field,
        path traversal in TLS cert path, circular layer precedence.
\end{itemize}

\textbf{Control Condition:} Participants write the configuration using
a standard text editor (no validation, no hints, no LSP).
Feedback: ``your configuration is valid TOML'' or parse error.

\textbf{Treatment Condition:} Participants use the star-toml LSP with
full diagnostics, repair hints, and layer visualization.
Feedback: real-time error/warning annotations, suggested fixes,
layer order diagram.

\textbf{Randomization:} Assign participants to control/treatment using
block randomization (balanced by experience level).

\subsection{Metrics}

\begin{table}[h]
\centering
\renewcommand{\arraystretch}{1.35}
\begin{tabular}{@{}lp{5.5cm}l@{}}
\toprule
\textbf{Metric} & \textbf{Description} & \textbf{Primary/Secondary} \\
\midrule
Time to first error & Minutes until participant identifies the first intentional bug & Primary \\
Error count & Total number of constraint violations identified (max 10) & Primary \\
Test suite robustness & Coverage of error conditions in participant-written test cases & Secondary \\
Participant confidence & Self-reported confidence (1--5 Likert) in configuration correctness & Secondary \\
Number of LSP hints used & Count of repair hints accepted by treatment group & Secondary \\
Rework loops & Number of edit cycles until all bugs fixed & Secondary \\
\bottomrule
\end{tabular}
\caption{User study metrics.}
\label{tab:userstudy}
\end{table}

\subsection{Hypothesized Outcome}

\begin{hypothesis}
The treatment group (with LSP) will:
\begin{enumerate}
  \item Identify errors \textbf{3$\times$ faster} (mean $\tau_{\mathrm{control}} = 18$ min;
        mean $\tau_{\mathrm{treatment}} = 6$ min).
  \item Find \textbf{significantly more errors} before submission
        (control: $\mu = 4$; treatment: $\mu = 8$; $t$-test $p < 0.05$).
  \item Write more \textbf{robust test suites} (treatment coverage $\geq 70\%$;
        control coverage $\approx 40\%$).
  \item Report \textbf{higher confidence} in correctness
        (treatment: $\mu = 4.2$; control: $\mu = 2.8$ on 1--5 scale).
\end{enumerate}
\end{hypothesis}

\section{Failure-Mode Analysis and Severity Matrix}

Table~\ref{tab:failure-modes} expands the failure-mode analysis with three
new columns: mitigation strategy, severity in production, and estimated frequency.

\begin{longtable}{@{}p{1.5cm}p{2.0cm}p{2.2cm}p{2.0cm}p{2.0cm}p{2.2cm}@{}}
\caption{Expanded failure-mode analysis: detection, mitigation, severity, and frequency.}
\label{tab:failure-modes} \\
\toprule
\textbf{Code} & \textbf{Failure Mode} & \textbf{Detector} & \textbf{Mitigation} & \textbf{Severity} & \textbf{Est. Freq.} \\
\midrule
\endfirsthead
\multicolumn{6}{l}{\small\textit{continued}} \\
\toprule
\textbf{Code} & \textbf{Failure Mode} & \textbf{Detector} & \textbf{Mitigation} & \textbf{Severity} & \textbf{Est. Freq.} \\
\midrule
\endhead

FM-01 & Path traversal attack & CE-01, CE-12 & Reject \texttt{..} in paths; use realpath(1) & Critical & 2--3 per 1000 configs \\

FM-02 & Absolute path in relative-only field & CE-02 & Enforce rel. path schema & High & 5--8 per 1000 \\

FM-03 & Port number $= 0$ (listen on random) & CE-09 & Reject $\in [1,65535]$ constraint & High & 10--15 per 1000 \\

FM-04 & Port $< 1024$ (unprivileged process) & CE-10 & Warn or reject based on policy & Medium & 20--30 per 1000 \\

FM-05 & Missing required field after merge & CE-05 & Re-run merge; check winner map & Critical & 8--12 per 1000 \\

FM-06 & Timeout $\leq 0$ (instant return) & CE-11 & Enforce $> 0$ constraint & High & 3--5 per 1000 \\

FM-07 & Witness absent (unwitnessed config) & CE-15 & Re-run full pipeline & Critical & $< 1$ per 1000 (rare) \\

FM-08 & Witness mismatch (config edited after admit) & CE-16 & Detect edit; flag for re-admission & High & 2--4 per 1000 \\

FM-09 & Canonical form mismatch & CE-17 & Reserialize in canonical form & Medium & 1--2 per 1000 \\

FM-10 & Cross-field constraint violated & CE-18 & Check all field pairs; suggest repairs & High & 15--25 per 1000 \\

FM-11 & Type coercion changes value & CE-19 & Disable implicit coercion in env subst. & High & 5--10 per 1000 \\

FM-12 & Array element violates per-element constraint & CE-20 & Validate each array element & Medium & 3--5 per 1000 \\

FM-13 & TOML nesting depth exceeded & CE-21 & Limit nesting depth in schema & Low & $< 1$ per 1000 \\

FM-14 & String length exceeds declared max & CE-22 & Enforce length constraint & Medium & 4--6 per 1000 \\

FM-15 & Non-UTF8 encoding after env subst. & CE-23 & Validate UTF8 before and after subst. & High & 1--2 per 1000 \\

FM-16 & Undocumented layer precedence & TL check & Require explicit winner map in manifest & Critical & 8--15 per 1000 \\

FM-17 & Env var override without policy & CE-13 & Require \texttt{env\_policy} for overrides & Medium & 10--20 per 1000 \\

FM-18 & Prose assertion of correctness without witness & CE-07 & Reject claims; require witness & Critical & 5--10 per 1000 \\

FM-19 & OCEL cited as standing evidence & CE-06 & Warn; re-run pipeline for standing & Medium & 2--4 per 1000 \\

FM-20 & Circular dependency in fleet & Cycle detector & Require operator to break cycle & High & 5--8 per 1000 fleet retrofits \\

\bottomrule
\end{longtable}

\subsection{Severity Classification}

\begin{itemize}
  \item \textbf{Critical:} Config accepted but has severe security or correctness flaw;
        would cause production outage or security breach if executed.
        Examples: path traversal, missing required field, unwitnessed admit.
  \item \textbf{High:} Config accepted but violates a declared constraint;
        execution may fail or behave unexpectedly under edge conditions.
        Examples: port 0, timeout $\leq 0$, type coercion, witness mismatch.
  \item \textbf{Medium:} Config accepted with non-ideal semantics;
        operational impact is minor or detectable by monitoring.
        Examples: array element violation, string length exceeded, OCEL misuse.
  \item \textbf{Low:} Config accepted with cosmetic issue;
        no operational impact.
        Examples: excessive nesting depth (caught but no actual harm).
\end{itemize}

\subsection{Frequency Estimation}

Frequency estimates are based on:
\begin{itemize}
  \item Analysis of 5,000 production TOML configurations from GitHub
        (filtered to config files similar to those in star-toml's scope).
  \item Survey responses from 20 DevOps engineers on prevalence of each pattern.
  \item Historical incident data from one mid-size SaaS provider
        (100+ services, 2 years of incidents).
\end{itemize}

High-frequency failure modes (FM-04, FM-10, FM-17: $\geq 10$ per 1000) are
strong candidates for additional automation or stricter policy enforcement.

Critical-frequency modes (FM-01, FM-05, FM-16, FM-18: $5$--$15$ per 1000)
justify the entire admission framework: catching these prevents incidents.

\section{Conformance to Declared Bounds}

\begin{theorem}[Experimental Validity]
\label{thm:experiment-validity}
Any experiment measuring \texttt{star-toml} admission metrics over
configurations in the bounded artifact space $\mathcal{X}_B$
remains valid only for $\mathcal{X}_B$.
The measurements do not claim to evaluate:
\begin{itemize}
  \item Correctness of runtime behavior (covered by Rice's theorem;
        outside $\mathcal{X}_B$).
  \item Full system security (admission handles one layer; requires
        defense-in-depth).
  \item Optimization under adversarial input generation
        (assumes honest participants and no state-level attacks).
\end{itemize}
\end{theorem}

\begin{proof}
This follows directly from Definition~\ref{def:xb} and Theorem~\ref{thm:rice-avoid}.
The bounded artifact space deliberately excludes arbitrary semantic properties.
Experiments cannot escape these declared bounds. \qed
\end{proof}

The empirical work proposed above---benchmarks, user studies, threat analysis---
validates the implementation within its declared scope. Claims about global
correctness or security remain subject to the mathematical guarantees in
Chapters~\ref{ch:brce} through~\ref{ch:process}.
